\section{Materials and methods}
\subsection{Cohort selection and independent units}
GSE165816 comprises 54 original samples and 27 author-mapped participants. The frozen model used 25 selected foot-skin samples in the non-diabetic, DFU-healing and DFU-non-healing groups: 81,602 cells before cell quality control and 76,461 afterwards. Cells required at least 200 detected genes and a mitochondrial count fraction at most 0.20; genes required detection in at least 10 cells before panel selection. The final panel contained 7,002 genes, obtained from a nominal 7,000-gene dispersion selection with retained anchor genes. The discovery outcome subset contained 9 healing and 5 non-healing samples.

The clinical supplement maps those 14 outcome samples to 11 patients: 7 healing and 4 non-healing. G4/G23, G7/G8 and G33/G34 are repeated samples from the same wound on the same day. We combine each patient's retained cells before computing the primary patient score. The outcome is the study's healing category within a 12-week follow-up window, not a measured individual healing time. Series metadata describe a shared Emory 10x Chromium 3$'$ V2/V3 workflow and NovaSeq instrumentation; the absence of per-sample chemistry and flow-cell identifiers prevents run-specific batch adjustment.

\subsection{What the frozen neural topic model computes}
For nonnegative panel counts $x_{ig}$ with positive total, let $u_{ig}=x_{ig}/\sum_{h=1}^{G}x_{ih}$. The implemented model uses a diagonal-Gaussian encoder in unconstrained coordinates and a softmax decoder:
\begin{align}
q_\phi(z\mid u_i)&=\mathcal N\!\left(\mu_\phi(u_i),\operatorname{diag}\sigma_\phi^2(u_i)\right),\label{eq:p1encoder}\\
\theta_{ik}&=\frac{e^{z_{ik}}}{\sum_{\ell=0}^{K-1}e^{z_{i\ell}}},\qquad
\beta_{kg}=\frac{e^{b_{kg}}}{\sum_{h=1}^{G}e^{b_{kh}}},\qquad
\widehat u_{ig}=\sum_{k=0}^{K-1}\theta_{ik}\beta_{kg}.\label{eq:p1decoder}
\end{align}
Equations~\eqref{eq:p1encoder}--\eqref{eq:p1decoder} define the encoder, simplex coordinates and normalized topic-gene decoder used below. Thus $\theta_i\in\Delta^{K-1}=\{a\in\R^K:a_k\geq0,\sum_{k=0}^{K-1}a_k=1\}$ and every topic-gene row sums to one. The topic indices start at zero. A Gaussian mapped through softmax gives a logistic-normal representation. This differs from a Dirichlet prior, while retaining the mixture representation that motivates topic models. Training minimizes
\begin{equation}
\mathcal L(\phi,b)=\frac1N\sum_i\E_{q_\phi}\!\left[-\sum_g u_{ig}\log(\widehat u_{ig}+10^{-10})\right]
+0.01\,\frac1N\sum_i\operatorname{KL}\!\left(q_\phi(z\mid u_i)\,\|\,\mathcal N(0,I)\right).\label{eq:p1loss}
\end{equation}
Equation~\eqref{eq:p1loss} is the fitted objective: each cell contributes cross-entropy on its panel proportions, followed by a Gaussian regularizer. It is neither a depth-weighted count likelihood nor an unmodified evidence lower bound for such a likelihood. With encoder log-variance $\ell_{ik}=\log\sigma_{ik}^2$, the reparameterization and regularizer are
\begin{equation}
z_{ik}=\mu_{ik}+e^{\ell_{ik}/2}\epsilon_{ik},\quad \epsilon_i\sim\mathcal N(0,I),\qquad
\operatorname{KL}_i=\frac12\sum_k\left(\mu_{ik}^2+e^{\ell_{ik}}-1-\ell_{ik}\right).\label{eq:p1gaussian}
\end{equation}
Equation~\eqref{eq:p1gaussian} is the diagonal-Gaussian calculation used by the implementation; one noise draw per cell estimates the reconstruction expectation in each training batch. The encoder has widths 256 and 128, batch normalization, ReLU activations and dropout 0.1. The recorded fit used $K=15$, 40,000 Adam steps, learning rate $10^{-3}$ and batches of 512. Deterministic projection disables encoder dropout and uses $z=\mu_\phi(u)$ and $\theta=\operatorname{softmax}(\mu_\phi(u))$, which is not generally $\E_q[\operatorname{softmax}(z)]$. Missing panel genes in external projections are set to zero. Normalization and quality control can alter biological interpretations, so this definition is kept distinct from other variance-stabilizing methods.

Topic-gene perplexity is summarized as $\exp\{-K^{-1}\sum_{k,g}\beta_{kg}\log\beta_{kg}\}$; it measures concentration of decoder weights, not superiority over other representations. The latent coordinate $\mu\in\R^K$ is not the simplex coordinate $\theta$. Softmax is unchanged by adding a common scalar to all coordinates of a cell; reconstruction therefore does not identify that offset. The Gaussian penalty constrains the fitted coordinates but does not give them a unique biological interpretation. Topic weights are decoder coefficients, not independently calibrated probabilities of membership in a biological cell type.

\subsection{Threshold fractions and the exact mean decomposition}
Let $Y_{si}=\theta_{si,\star}$ for fibroblast-gated cell $i$ in sample $s$. The discovery lineage gate averages marker expression standardized across the selected discovery cells, then assigns the highest-scoring lineage. Its markers are disjoint from the tested topic-gene list by an explicit overlap check. The gate remains an expression-based operational definition; its reference population can change its assignments, and the frozen encoder alone does not freeze the complete tissue measurement. A pooled two-Gaussian fit provides a descriptive cut $c=0.1842958$, the midpoint of the fitted component means. Gaussian components approximate a bounded score distribution and are not assumed to be literal biological populations.
\begin{equation}
\bar Y_s=\frac1{n_s}\sum_iY_{si},\qquad
\widehat w_s(c)=\frac1{n_s}\sum_i\ind\{Y_{si}>c\}.\label{eq:p1fraction}
\end{equation}
Equation~\eqref{eq:p1fraction} defines both the sample mean and the reported hard-threshold fraction; the latter is not an average posterior responsibility of the Gaussian mixture. Define $\bar Y_{s,L}$ and $\bar Y_{s,H}$ as means below/at and above the cut. Then, with empty-component contributions interpreted as zero,
\begin{equation}
\bar Y_s=(1-\widehat w_s)\bar Y_{s,L}+\widehat w_s\bar Y_{s,H}.\label{eq:p1identity}
\end{equation}
Equation~\eqref{eq:p1identity} is exact. Only under an additional approximate component-stability assumption, $\bar Y_{s,L}\approx m_L$ and $\bar Y_{s,H}\approx m_H$, does it reduce to $\bar Y_s\approx m_L+(m_H-m_L)\widehat w_s$. Consequently, a large correlation between the two summaries cannot independently prove stable intensity: both summaries use the same cells. Persistence of overlapping low-state tails supplies complementary distributional evidence. The 10th percentile of a high-loading sample (0.029) overlaps the scale of the 90th percentile in a low-loading sample (0.119); this observation alone is not a formal test of multimodality.

To quantify the term omitted by a constant-state approximation, let $m_L,m_H$ be the cell-pooled means in the two threshold bins across the 19,410 discovery fibroblasts. They are empirical bin means, not the fitted Gaussian means. For each specimen,
\begin{align}
A_s&=(1-\widehat w_s)m_L+\widehat w_sm_H,\notag\\
R_s&=(1-\widehat w_s)(\bar Y_{s,L}-m_L)+\widehat w_s(\bar Y_{s,H}-m_H),\qquad \bar Y_s=A_s+R_s.\label{eq:p1residual}
\end{align}
Equation~\eqref{eq:p1residual} separates a specified constant-state approximation from its exact residual. An empty bin contributes zero and its undefined mean is omitted from the corresponding mean range. We summarize all 25 specimens using specimen-weighted residual RMS and bin-mean ranges. These are finite-data descriptions without a test of state invariance or a causal allocation of variance (Table~\ref{tab:p1math}).

The semantic statistic is $\rho_{\rm sem}=\operatorname{cor}_{\rm Spearman}(\bar Y_s,\widehat w_s)$ across 25 samples. A 10,000-draw sample-row percentile bootstrap, 10,000 row permutations and leave-one-sample-out recalculation characterize its stability. These do not create independent patients. The patient-collapsed correlation of 0.937 across 20 mapped patients and its descriptive patient-resampling interval are the primary uncertainty summary; the earlier specimen resampling is historical.

A separate patient-resampling analysis collapses the 25 specimen means and high-state fractions to 20 patients using fibroblast counts as within-patient weights. It draws 10,000 patient samples with replacement (seed 217) and reports the percentile interval for their Spearman correlation. Leave-one-patient-out analyses recompute this correlation across the remaining 19 patients and separately recompute the all-cell healing contrast across the remaining 10 outcome patients. The reported omission range is a sensitivity diagnostic, not an interval estimate (Table~\ref{tab:p1robustness}).

\subsection{Technical controls and constant sensitivity}
Dissociation-induced programmes and ambient transcripts motivate the technical controls. The stress comparison uses a standardized mean difference (SMD) with rejection bound 0.8; the doublet-oriented control compares PTPRC-positive rates with ratio bound 2.0. Exact panel-count matching reprojects downsampled counts, comparing cell labels and sample fractions. The historical 31-covariate screen retains its specimen-row asymptotic p values and BH transforms for auditability, not as valid patient-level significance or false-discovery-rate control: repeated patients violate the naive independent-row premise. Removing any cell carrying a specified immune transcript tests invariance to that exclusion rule. This cannot by itself identify ambient contamination or exclude every doublet. Residual mixture checks with immune covariates are compared with four unrelated placebo covariates; similarly flattened residuals limit claims of statistical separability.

Sensitivity fits use 18 cuts between 0.08 and 0.40 (including the rounded fitted cut) and $K\in\{10,15,20\}$. The TIMP1-carrying topic is tracked descriptively across capacities; topic labels are not assumed exchangeable or biologically identical across $K$. The anatomical positive-control calibration uses 20,000 shuffles. The raw-count structure control selects a fixed 6,000-cell subset by the recorded source order, then permutes each gene's observed values across cells with a seeded permutation; this preserves the marginal multiset for each gene but does not preserve row-wise library sizes. The resulting topic concentration, row-wise library-size and detection-rate summaries are compared with real input; no new topology audit is performed. The separate outcome-label calibration exhaustively enumerates all 2,002 allocations of the 9 healing and 5 non-healing specimens, retaining repeated specimens as a specimen-unit null. It is not a patient-level clinical test or proof of pipeline-wide error control. The DFU-only label null excludes all healthy controls before allocation and reports the 9/5 specimen result as secondary to the 7/4 patient test. For each specimen, a continuous-score construction applies the inverse-logit transform to its observed mean on the logit scale plus independent standard-normal noise, retaining its cell count; 200 seeded draws are made without a latent class or mixture fit. The constructed mean need not equal the observed mean. This is an illustrative non-identification example, not a fitted null model or a calibrated probability. These analyses are retrospective computational controls, not a registered clinical trial.

For marker-reference sensitivity, expression is normalized per cell to 10,000 counts and transformed by log1p. Each lineage score averages its marker-wise standardized values, followed by exclusive maximum-score assignment. The pooled marker means and SDs are frozen before applying the same reference to each specimen; this implementation check is compared with re-estimating means and SDs separately within each specimen. The topic encoder and state threshold remain fixed. Patient contrasts are recomputed after count-weighted pooling under each annotation reference, with the three readouts forming a separate BH family within each gate. This evaluates dependence on the annotation reference without an independent cell-identity truth set (Table~\ref{tab:p1gatecontrol}).

The exploratory covariate scan retains the original rounded rule $Y_{si}\geq0.184$, used by the technical-control analyses, rather than substituting the fitted cut 0.1842958. It correlates the sample high-state fraction with 14 other fibroblast topic means, 11 whole-sample lineage fractions, total and fibroblast cell counts, fibroblast dissociation score, median UMI count, PTPRC-positive fraction and the numeric G specimen label. Two-sided Spearman probabilities are adjusted across all 31 comparisons by Benjamini--Hochberg. The 25 rows are specimens, include repeated patients, and receive no clinical-outcome adjustment. The G label is an identifier covariate, not a verified sequencing batch or collection time. Topic coordinates and lineage fractions have compositional dependence.

\subsection{Full released-count depth sensitivity}
The full-pipeline stress test starts from all 81,602 released count columns in the 25 selected specimens, including cells failing the original analysis QC. These are processed count matrices, not pre-cell-calling droplets or sequencing reads. Baseline reconstruction must reproduce retained cell identities, lineage assignments and deterministic coordinates before perturbation. Each nonzero count is independently thinned with a binomial probability of 0.75, 0.50, 0.25 or 0.10, with three fixed computational seeds per rate. Every rate starts from the original counts; no cells or patients are resampled. QC is rerun using at least 200 detected genes and a mitochondrial fraction at most 0.20. The discovery marker reference, exclusive maximum-score gate, panel, encoder and high-state threshold remain unchanged; there is no new margin or outcome-guided label rule.

The full-pipeline population includes all current QC survivors and uses newly computed labels. Separate conditional branches retain original QC cells and labels, retain the QC intersection with original labels, or re-gate the same intersection. Lost and newly admitted cells, lineage transitions and fibroblast entry/exit are reported separately. Zero detected annotation markers and tied winning scores are diagnostics, not evidence of valid cell identity. A zero panel total makes a score missing, rather than assigning the encoder output of a zero vector. Readouts preserve their selected-cell denominators and are non-estimable if required scores are missing. Patient aggregation pools the appropriate all-cell or fibroblast counts across repeated specimens. The same seven healed and four non-healed patients supply 330 exhaustive label allocations; the three seed ranges describe perturbation sensitivity, not biological confidence intervals or new outcome tests. Uniform thinning does not model upstream cell calling, cell-type-specific capture, ambient RNA, doublets or independent tissue sampling.

\subsection{Patient contrast, permutation and equivalence sensitivity}
Let $\mathcal C_p$ contain all retained cells in the mapped DFU foot specimens of patient $p$. The primary corrected score is $S_p=|\mathcal C_p|^{-1}\sum_{i\in\mathcal C_p}\theta_{i\star}$, so repeated specimens contribute in proportion to retained cell counts. This all-cell readout differs from $\bar Y_s$ restricted to fibroblasts and from $\widehat w_s$. For healing set $H$ and non-healing set $N$,
\begin{equation}
\widehat\Delta=\frac1{|H|}\sum_{p\in H}S_p-\frac1{|N|}\sum_{p\in N}S_p.\label{eq:p1delta}
\end{equation}
Equation~\eqref{eq:p1delta} is the primary patient-unit contrast; it is separate from the sample-level fibroblast ratio and from the threshold fraction. All $\binom{11}{7}=330$ patient-label allocations are enumerated. To reproduce the recorded result, the implementation retains an add-one correction even for exhaustive enumeration:
\begin{equation}
p_{+}=\frac{1+\sum_{\pi\in\Pi}\ind\{|\widehat\Delta_\pi|\geq|\widehat\Delta|\}}{1+|\Pi|}.\label{eq:p1permutation}
\end{equation}
Equation~\eqref{eq:p1permutation} is the recorded add-one-corrected two-sided permutation probability. It is slightly conservative relative to the uncorrected exhaustive fraction; the add-one convention is usually motivated for randomly drawn permutations, whereas all allocations are enumerated here. Exchangeability is an assumption for an observational label comparison, not a randomization supplied by study design. A 30,000-draw bootstrap resamples patients within outcome arms. Small group sizes and conditional resampling explain why bootstrap intervals and permutation tests need not be dual.

An independent enumeration also reports the exhaustive probability without the extra correction:
\begin{equation}
p_{\rm exact}=\frac{\#\{\pi\in\Pi:|\widehat\Delta_\pi|\geq|\widehat\Delta|\}}{|\Pi|}=\frac{90}{330}=0.272727.\label{eq:p1exact}
\end{equation}
Equation~\eqref{eq:p1exact} includes the observed allocation and numerical ties. It distinguishes exhaustive enumeration from randomly sampled permutations; the recorded add-one comparison $91/331=0.274924$ is retained alongside it. Neither convention makes observational healing labels randomized or tests a causal effect. Table~\ref{tab:p1exact} reports the independent arithmetic check, including the Welch quantities below.

For a symmetric equivalence margin $M$, the Welch two-one-sided-test procedure uses $SE=(s_H^2/n_H+s_N^2/n_N)^{1/2}$ and Welch degrees of freedom $\nu$:
\begin{equation}
p_L=1-F_{t_\nu}\!\left(\frac{\widehat\Delta+M}{SE}\right),\qquad
p_U=F_{t_\nu}\!\left(\frac{\widehat\Delta-M}{SE}\right).\label{eq:p1tost}
\end{equation}
Equation~\eqref{eq:p1tost} gives the two one-sided tests used for a prespecified margin. Both must be below 0.05 to reject non-equivalence. Writing $v_H=s_H^2/n_H$ and $v_N=s_N^2/n_N$ makes the degrees of freedom and margin boundary explicit:
\begin{equation}
\nu=\frac{(v_H+v_N)^2}{v_H^2/(n_H-1)+v_N^2/(n_N-1)},\qquad
M_*=|\widehat\Delta|+t_{\nu,0.95}SE.\label{eq:p1margin}
\end{equation}
Equation~\eqref{eq:p1margin} is the infimum margin containing the two-sided 90\% Welch interval; strict rejection requires $M>M_*$. Independent calculation gives $SE=0.019446$, $\nu=8.113833$ and $M_*=0.065375$ for the 11-patient comparison. The deterministic specimen-unit boundary is 0.079644. A margin selected after seeing these data is a sensitivity summary, not a clinically justified equivalence claim. No clinical minimum important difference has been supplied.

\subsection{Discrimination, anatomy and design simulations}
For held-out prediction scores $r$, AUC is the empirical probability that a healing score exceeds a non-healing score, with half credit for ties:
\begin{equation}
\widehat{\rm AUC}=\frac1{n_Hn_N}\sum_{h\in H,n\in N}\left[\ind\{r_h>r_n\}+\tfrac12\ind\{r_h=r_n\}\right].\label{eq:p1auc}
\end{equation}
Equation~\eqref{eq:p1auc} is the pairwise empirical AUC used for the held-out scores. Topic orientation, module preprocessing and the choice of PCA/NMF components are determined inside each training fold. The topic encoder itself was trained without outcomes on the discovery cohort, including held-out patients' expression: this is an internal discrimination audit of a frozen discovery representation, not fully external representation validation. Bootstrap intervals resample fixed out-of-fold scores with common arm-stratified indices across methods; the entire fitting/orientation process is rerun for label permutations. All allocations are used for topic and module readouts; PCA and NMF use 200 sampled relabellings. Sample-level comparisons use 14 specimens and patient-level comparisons use 11 mapped patients; they are not interchangeable.

PCA scaling and basis fitting, and NMF fitting, use training rows only. The historical pooled out-of-fold AUC nevertheless compares scores from different fitted folds; those scores are not calibrated clinical risks. Separate PCA and NMF fits can change the selected coordinate, origin and scale between folds. AUC is invariant to a common increasing transformation, but not to separate fold-specific transformations. These raw-coordinate AUCs and their fixed-score bootstrap contrasts are retained as historical scoring diagnostics, not evidence for a representation ranking.

The revised patient discrimination audit instead leaves out one healing patient $h$ and one non-healing patient $n$ together. Each of the $7\times4=28$ folds fits one scoring function $r^{(-h,-n)}$ to the remaining nine patients (six healing, three non-healing), and that unchanged function scores both excluded patients:
\begin{equation}
\widehat A_{\rm pair}=\frac1{28}\sum_{h\in H,n\in N}\left[\ind\{r^{(-h,-n)}_h>r^{(-h,-n)}_n\}+\tfrac12\ind\{r^{(-h,-n)}_h=r^{(-h,-n)}_n\}\right].\label{eq:p1pairauc}
\end{equation}
Equation~\eqref{eq:p1pairauc} averages within-model pair orderings; no raw scores from different fitted models are compared. Patient expression inputs are cell-count-weighted means of per-cell proportions on the fixed panel. Topic orientation, module gene standardization and orientation, PCA standardization and fitting, and NMF fitting and component selection use training patients only. The PCA rule selects the largest absolute training-arm mean difference among five components. NMF uses five components and tolerance $10^{-4}$; the initial 1,000-iteration cap was raised uniformly to 10,000 after three training fits reached the cap, without using held-pair discrimination to select it. The initial run is retained, and no final fit emitted a convergence warning. Fitted basis rows are normalized to unit Euclidean norm before independent nonnegative-least-squares projection of each row. The component is selected by the absolute training-arm mean difference divided by the training-patient SD, making selection invariant to positive component rescaling. The direction is set by the training contrast, and exact ties receive half credit. The fixed gene panel, topic representation and study-specific analysis choices remain discovery-derived. The 28 pairs share 11 patients and overlapping training sets, so they are not 28 independent outcomes; no pair-binomial uncertainty, fixed-score bootstrap interval or significance probability is supplied for this revised estimand. Changing component normalization and selection as well as fold structure means the difference from the historical AUC cannot identify the contribution of one correction alone. Reusing the same cohort to inspect readouts limits confirmatory interpretation.

The paired anatomical control tests foot-minus-forearm differences in 10 mapped patients using all $2^{10}$ sign allocations and the same add-one convention. Across 15 topics, Benjamini--Hochberg adjustment is $q_{(i)}=\min\{1,\min_{j\geq i}15p_{(j)}/j\}$, with the usual assumptions required for false-discovery-rate interpretation.

The historical high-state-fraction power calculation uses 9 versus 5 specimens, Gaussian location shifts with pooled SD 0.283 and a two-sided Mann--Whitney test. Its 24.9\% power and $\operatorname{MDE}_{0.8}=0.519$ apply to this approximate design, not to the 11-patient all-cell contrast. Gaussian simulation can leave the unit interval, so it is a design-scale sensitivity approximation. A separate next-cohort simulation uses specimen-derived arm SDs of deterministic all-cell mean scores, 0.05170 and 0.01686, Welch testing or Welch equivalence testing, and 4,000 draws per candidate size. These descriptive spreads include repeated specimens and are not independent-patient variance estimates. The candidate totals are approximate conditional calculations, not validated recruitment targets (Table~\ref{tab:p1design}).

\subsection{External projections and limits of identification}
Every external cohort is projected into the frozen panel/model without outcome-guided retuning. The immune and B/plasma gates exclude promiscuous markers such as CD74, CD37, CD38 and SDC1 and assign non-overlapping compartments. Immune fraction uses all retained cells as its denominator; B/plasma share within immune cells uses only that compartment. Their correlations estimate different quantities. Negative depth correlations cannot by themselves exclude depth confounding, because a marginal sign is not a conditional adjustment.

GSE231643 supplies eight title-labelled specimens but lacks an authoritative patient mapping. GSE241132 contributes an acute-wound, three-donor construct summary. GSE223964, GSE245703 and GSE268834 lack qualifying patient-level healing endpoints, and GSE248247 has only two samples. We do not pool these correlations or describe them as independent prognostic validation. No analysis identifies a treatment effect, a time-to-healing model or a causal effect of immune infiltration.

\subsection{Observed-cell and expression displays}
The biological displays reuse the saved 15-topic cell proportions and operational lineage assignments. For visualization only, each proportion vector is transformed to centered log ratios,
\begin{equation}
c_{ik}=\log\theta_{ik}-K^{-1}\sum_{\ell=0}^{K-1}\log\theta_{i\ell}.\label{eq:p1display}
\end{equation}
Equation~\eqref{eq:p1display} removes the geometric mean of the positive coordinates before a Euclidean-distance UMAP with 30 neighbors, minimum distance 0.3 and random seed 17. All 76,461 cells enter the embedding; a uniform subset of 2,000 cells is drawn for the printed vector panels, using the same cells and coordinates for lineage and loading colors. The sampling affects only display. Neither embedding separation nor displayed point density enters the outcome tests, gate, threshold or model fitting. Two-dimensional maps can distort relationships in the input space and do not establish cell types or trajectories.

Raw discovery matrices are reloaded under the original tissue, disease and cell-quality criteria. We require the retained sample and arm order and all recomputed marker-based labels to match the saved cell table. For each displayed gene, dot size is the fraction of cells with a positive count; color is the mean of cellwise $\log(1+10^4 x_{ig}/\sum_h x_{ih})$, scaled from its smallest to largest lineage mean. The denominator uses the source gene matrix, not only the frozen modeling panel. The 18-gene display includes compartment markers and programme-associated genes without an outcome-based selection step. It illustrates measured expression and cannot independently validate labels derived from some of the same markers. Specimen composition retains each specimen as a separate column, and loading histograms use all assigned fibroblasts in fixed bins of width 0.02.

\subsection{Patient-unit extension of the lineage readouts}
The retrospective extension fixes three readouts before computing their new outcome comparisons: within-fibroblast mean loading, the fraction above the existing mixture midpoint, and fibroblast composition among all retained cells. Let $\mathcal S_p$ denote patient $p$'s specimens, $n_s^F$ their fibroblast counts and $n_s^A$ their all-cell counts. With specimen means $m_s$, threshold fractions $w_s$ and composition fractions $f_s=n_s^F/n_s^A$, the pooled patient summaries are
\begin{equation}
M_p=\frac{\sum_{s\in\mathcal S_p}n_s^F m_s}{\sum_{s\in\mathcal S_p}n_s^F},\quad
W_p=\frac{\sum_{s\in\mathcal S_p}n_s^F w_s}{\sum_{s\in\mathcal S_p}n_s^F},\quad
F_p=\frac{\sum_{s\in\mathcal S_p}n_s^A f_s}{\sum_{s\in\mathcal S_p}n_s^A}.\label{eq:p1lineagepatients}
\end{equation}
Equation~\eqref{eq:p1lineagepatients} pools the appropriate denominators within patient while giving each of the 11 outcome patients equal weight in the group contrast. The representation, annotations and loading cut are unchanged. We enumerate all 330 allocations of seven healed labels among 11 patient rows, counting absolute contrast ties within a numerical tolerance of $10^{-12}$. For this new exhaustive enumeration the probability is the exceeding count divided by 330, without the add-one convention retained in the historical all-cell analysis. The three new probabilities form one Benjamini--Hochberg family. Observational group exchangeability is an assumption, not random assignment of healing outcomes.

Percentile intervals use 10,000 patient resamples separately within the seven healed and four non-healed groups, with seed 17. They are small-sample descriptive intervals, not externally calibrated prediction intervals. Sensitivities replace within-patient count weights by equal specimen weights, omit each patient in turn, and require every retained specimen to have at least 100 fibroblasts. Common random draws are used when the patient arrays coincide. We report when an eligibility sensitivity removes no specimen. The 20-patient descriptive table includes non-diabetic participants, but only the 11 DFU patients enter healing contrasts; no healthy participant is assigned a healing outcome.

\subsection{Patient distributions, cell yield and repeated sampling}
To assess the loading distribution without selecting one state threshold, let $F_p(c)$ be the empirical cumulative distribution of all recovered fibroblast loadings pooled within patient $p$. Each patient contributes a normalized distribution regardless of cell yield. We compare the two equally weighted patient-group distributions using
\begin{equation}
T=\sup_{c\in\mathcal C}\left|\frac1{7}\sum_{p:H_p=1}F_p(c)-\frac1{4}\sum_{p:H_p=0}F_p(c)\right|,\label{eq:p1ecdf}
\end{equation}
where $\mathcal C$ contains every distinct observed loading. Equation~\eqref{eq:p1ecdf} uses the same 11 patients and frozen scores. Its exact probability enumerates all 330 patient-label allocations; this is a patient-level permutation statistic, not a cell-level Kolmogorov--Smirnov test. We also report high-fraction contrasts at cuts 0.05, 0.10, 0.15, 0.20, 0.30, 0.40, 0.50 and 0.60, together with the original midpoint. Comparing each absolute contrast with the permutation distribution of $T$ gives a max-statistic probability across the complete loading support. These retrospective sensitivities and their unadjusted probabilities are shown in Table~\ref{tab:p1thresholdpatient}.

Cell-yield sensitivity uses 250 draws of exactly 100 fibroblasts without replacement from each of the 14 DFU specimens, with random seed 29. Each draw averages specimen means and high-state fractions equally within patient before recomputing the 7-versus-4 contrast and exact probability. This changes the contribution of recovered cells, without changing clinical labels or the frozen threshold. Fibroblast composition is not recomputed from this deliberately selected cell subset. Separately, three patients have two specimens each; all $2^3=8$ choices of one specimen per patient are evaluated for all three original readouts. Table~\ref{tab:p1samplingsensitivity} gives complete contrast and probability ranges. The draws and specimen choices assess finite-cell and sampling sensitivity, not new patients or confidence bounds.


Repeated-specimen diagnostics retain all five mapped pairs, including three DFU pairs. For each pair we compare fibroblast means, high-state fractions and the supremum difference between empirical loading distributions. Within each specimen, 250 seeded random partitions create disjoint fibroblast halves without replacement; their absolute mean differences provide a conditional cell-sampling reference. We compare each observed between-specimen difference with both within-specimen 97.5th percentiles. These partitions are not biological replicates, patient confidence intervals or an intraclass correlation estimate (Table~\ref{tab:p1repeats}).
